% Notas de aula — Capítulo 19: Dispersão de Poluentes
% Curso: Meteorologia Prática (baseado em R. Stull, Practical Meteorology, CC BY-NC-SA 4.0)
% Caixas verdes = conteúdo literal do quadro; linhas verdes = encenação.
\documentclass[11pt]{article}
\usepackage{../comum/preambulo}
\graphicspath{{../figuras/cap19/}}

\title{Capítulo 19 --- Dispersão de Poluentes\\
{\large Notas de aula (roteiro de quadro-negro)}}
\author{Curso de Meteorologia Prática}
\date{}

\begin{document}
\maketitle
\thispagestyle{fancy}

\noindent\textbf{Plano sugerido:} 2 aulas de 2\,h.
\emph{Aula 1:} \S\ref{sec:pluma}--\S\ref{sec:estab} (pluma gaussiana,
imagem no chão, estabilidade e formas de pluma).
\emph{Aula 2:} \S\ref{sec:caixa}--\S\ref{sec:walk} (modelo de caixa,
episódios, passeio aleatório e a difusão turbulenta).

\section{O problema}\label{sec:pluma}

\begin{noquadro}[abertura]
\textbf{Cap.~19 --- Dispersão de Poluentes}\\[2pt]
a pergunta regulatória: dada uma fonte $Q$, qual a concentração $C$
em cada lugar?\\
a resposta é 100\% meteorologia: vento ($U$), turbulência
($\sigma$s) e teto ($z_i$)\\[2pt]
``a solução da poluição não é a diluição — mas a diluição é a
física''
\end{noquadro}

Este capítulo é o Cap.~18 aplicado: a CLA \javimos{18}{camada
limite} vira o recipiente, a turbulência \javimos{18}{TKE e espectro}
vira o agente difusor, e as estabilidades do Cap.~5
\javimos{5}{classes de estabilidade} esculpem as plumas que qualquer
um vê de uma chaminé.

\section{A pluma gaussiana}\label{sec:gauss}

\quadro{Escrever a equação de advecção-difusão, trocar $t = x/U$ e
mostrar que é a equação do calor: a gaussiana não é chute, é
solução.}

\begin{formadif}[dedução (T1) e imagem (T2)]
\[ C(x,y,z) = \frac{Q}{2\pi U\sigma_y\sigma_z}\,
e^{-y^2/2\sigma_y^2}\left[e^{-(z-H)^2/2\sigma_z^2} +
e^{-(z+H)^2/2\sigma_z^2}\right], \]
com $\sigma^2 = 2Kx/U$. O segundo termo é a \textbf{fonte-imagem} em
$z = -H$: o método das imagens da eletrostática garantindo fluxo nulo
no chão — no solo, $C$ \emph{dobra}.
\end{formadif}

\begin{noquadro}[o máximo no solo (T3)]
perto da chaminé: $C_{solo} \approx 0$ (a pluma ainda não desceu)\\
máximo em $\sigma_z = H/\sqrt2$ \quad e \quad
$C_{max} \propto 1/H^2$\\[2pt]
dobrar a chaminé $=$ 1/4 do pico — mas cai mais LONGE (N1): a
política das chaminés altas exportou chuva ácida
\end{noquadro}

\section{Estabilidade esculpe a pluma}\label{sec:estab}

\quadro{Desenhar as três plumas clássicas (looping, coning, fanning)
e pedir à turma que diga a hora do dia e o perfil de $\theta_v$ de
cada uma \javimos{5}{perfis}.}

\begin{itemize}
  \item \textbf{Looping} (tarde instável): térmicas jogam a pluma no
        chão em golfadas — picos altos e curtos.
  \item \textbf{Coning} (neutro, vento forte): o cone gaussiano de
        livro.
  \item \textbf{Fanning} (noite estável): fita fina que viaja
        quilômetros sem descer\ldots
  \item \ldots até a \textbf{fumigação} da manhã: a camada de mistura
        crescente \javimos{18}{crescimento de $z_i$} alcança a fita e
        despeja tudo no chão de uma vez (N2) — o pico matinal das
        cidades industriais.
\end{itemize}

As classes de Pasquill (A--F) quantificam: $\sigma_y(x), \sigma_z(x)$
por estabilidade — a versão tabelada do Caso~2 do notebook.

\section{O modelo de caixa e os episódios}\label{sec:caixa}

\begin{formadif}[a caixa da cidade (T4; Caso~3)]
\[ z_i\frac{dC}{dt} = Q_A - \frac{U\,z_i}{W}\,C
\quad\Rightarrow\quad
C_{eq} = \frac{Q_A\,W}{U\,z_i}. \]
$U z_i$ $=$ \textbf{fator de ventilação}: a vazão de ar limpo. Inverno
estagnado (inversão baixa \javimos{5}{inversões} $+$ calmaria
anticiclônica \javimos{12}{altas de inverno}): $Uz_i$ despenca e a
mesma emissão dá 3--4$\times$ mais concentração — os episódios de
junho--agosto de São Paulo são meteorologia, não pico de emissão.
\end{formadif}

O N3 acrescenta o ciclo do tráfego e explica o clássico máximo
noturno: o rush das 18\,h emite dentro de uma caixa que está
\emph{encolhendo} ($z_i$ caindo \javimos{18}{transição noturna}).

\section{O passeio aleatório e a difusão turbulenta}\label{sec:walk}

\quadro{Fechar o curso de dispersão com Einstein: cada bocado de
fumaça é uma partícula browniana com a turbulência no papel das
moléculas.}

\begin{formadif}[Langevin e Taylor (T5; Caso~4)]
Velocidade com memória $T_L$ (escala lagrangiana):
\[ \sigma_z^2(t) =
\begin{cases}
\sigma_w^2\,t^2 & t \ll T_L\ \text{(balístico)}\\
2K\,t,\quad K = \sigma_w^2T_L & t \gg T_L\ \text{(difusivo)}
\end{cases} \]
O cotovelo de Taylor (1921) aparece limpo na simulação de partículas
do Caso~4 — e os modelos regulatórios (AERMOD, CALPUFF) são
exatamente esse passeio aleatório em escala industrial.
\end{formadif}

Deposição, química (ozônio fotoquímico, chuva ácida) e ilha de calor
urbana fecham o quadro qualitativamente — cada um mereceria capítulo
próprio; o essencial: a meteorologia controla o \emph{quando} e o
\emph{quanto} de todos eles \veremos{21}{química e clima}.

\section{Síntese da aula}

\begin{noquadro}[mapa final --- canto do quadro]
\[ C = \frac{Q}{2\pi U\sigma_y\sigma_z}e^{\cdots}\ (+\text{imagem})
\qquad C_{max} \propto \frac{1}{H^2} \qquad
C_{eq} = \frac{Q_AW}{Uz_i} \qquad K = \sigma_w^2T_L \]
a poluição de uma cidade é o Cap.~18 com uma fonte dentro
\end{noquadro}

\section{Exercícios complementares}\label{sec:exercicios}

Soluções (professor) em \texttt{cap19\_solucoes.ipynb}.

\subsection*{Teóricos}

\begin{enumerate}[label=\textbf{T\arabic*.},itemsep=4pt]
  \item Resolva a advecção-difusão estacionária com $t = x/U$ e
        obtenha a pluma gaussiana com
        $\sigma^2 = 2Kx/U$; explique o prefator pela conservação de
        massa.
  \item Justifique a fonte-imagem pela condição de fluxo nulo no chão
        e mostre que a concentração no solo dobra; descreva a
        generalização para a tampa em $z_i$.
  \item Mostre que o máximo de concentração no solo ocorre onde
        $\sigma_z = H/\sqrt2$ e que $C_{max} \propto 1/H^2$.
  \item Deduza o modelo de caixa e o equilíbrio
        $C_{eq} = Q_AW/(Uz_i)$; defina o fator de ventilação e
        calcule-o para um dia ventilado e um episódio.
  \item Derive $\sigma^2 = 2Kt$ do passeio aleatório e conecte com o
        teorema de Taylor (regimes balístico e difusivo).
\end{enumerate}

\subsection*{Numéricos (no notebook)}

\begin{enumerate}[label=\textbf{N\arabic*.},itemsep=4pt]
  \item Varra a altura da chaminé e verifique
        $C_{max}\propto 1/H^2$ e o deslocamento de $x_{max}$.
  \item Simule a fumigação da manhã com a pluma noturna capturada
        pela CL crescente.
  \item Imponha o ciclo do tráfego na caixa e explique o máximo
        noturno de concentração.
  \item Varra $T_L$ no passeio aleatório e confirme
        $K = \sigma_w^2T_L$; identifique o regime balístico perto da
        fonte.
\end{enumerate}

\vspace{1em}
\noindent\rule{\linewidth}{0.4pt}\\
{\scriptsize Material derivado de R.~Stull, \emph{Practical Meteorology}
(CC BY-NC-SA 4.0); este documento herda a mesma licença.}

\end{document}
